\subsection{მიღებული შედეგი და შეფასება}

პროექტის საბოლოო შედეგი არის მომუშავე, კროსპლატფორმული desktop
აპლიკაცია, რომელიც Linux-ზე, Windows-სა და macOS-ზე ეშვება. ტიპური
workflow ასე გამოიყურება: პირველი გაშვებისას მომხმარებელი სახელს
ირჩევს (იმახსოვრება შემდეგი გაშვებებისთვის), შემდეგ ხსნის ან ქმნის
დოკუმენტს და გვერდითი პანელიდან იწყებს სესიას (ფაილების სექცია, იხ.
სურათი~\ref{fig:files}). აპლიკაცია რამდენიმე წამში
აბრუნებს ორ მოსაწვევ ბმულს, ლოკალური ქსელისა და საჯაროს (იხ.
სურათი~\ref{fig:host}). სტუმრები საკუთარ აპლიკაციაში ბმულს ჩასვამენ
და მაშინვე ხედავენ დოკუმენტის მიმდინარე მდგომარეობას; ყველა შემდგომი
ცვლილება ყველასთან რეალურ დროში აისახება. ჰოსტი peers პანელიდან
მართავს, ვის შეუძლია რედაქტირება (იხ. სურათი~\ref{fig:guest}), ხოლო
ნამუშევარი ნებისმიერ დროს ინახება როგორც ჩვეულებრივ ტექსტურ ფაილად, ისე
\texttt{.pcdoc} ფორმატში (სრული ისტორიით). სესიის მდგომარეობის
სადიაგნოსტიკოდ ჩაშენებულია მეტრიკების პანელი (იხ.
სურათი~\ref{fig:metrics}).

\appscreenshot{images/1.png}{ფაილების სექცია: შენახვა, გახსნა, Fork და
ბოლო ფაილების სია}{fig:files}

\appscreenshot{images/2.png}{ჰოსტის ხედი: gateway/tunnel სტატუსი,
მოსაწვევი ბმულები და Invite ღილაკი (statusline-ზე ჩანს HOST როლი და
room id)}{fig:host}

\appscreenshot{images/3.png}{სტუმრის ხედი read-only რეჟიმში და peers
პანელი, სადაც ჩანს ჰოსტი და read-only სტუმარი}{fig:guest}

\appscreenshot{images/4.png}{მეტრიკების პანელი: გეითვეის live
სტატისტიკა (კლიენტები, ოთახები, გადაგზავნილი ფრეიმები, uptime)}{fig:metrics}

\paragraph{მიზნების შესრულება.}
შესავალში ჩამოყალიბებული ოთხივე მიზანი მიღწეულია: (1) სესიის დაწყება და
ბმულით შეერთება მუშაობს ორივე ქსელურ რეჟიმში, ლოკალურ ქსელში
ინტერნეტის გარეშე და გლობალურად Cloudflare Tunnel-ით; (2) პარალელური
რედაქტირება უკონფლიქტოდ ერწყმის YATA CRDT-ით და მონაწილეები
გარანტირებულად ერთსა და იმავე ტექსტს ხედავენ; (3) ინფრასტრუქტურის მართვა
სრულად ავტომატიზებულია, მომხმარებელი სერვერს ვერც ამჩნევს; (4) ჰოსტი
სესიას სრულად აკონტროლებს უფლებების ინდივიდუალური მართვით.

\paragraph{ხარისხის უზრუნველყოფა.}
სამივე კომპონენტს საკუთარი ტესტები აქვს: CRDT ბიბლიოთეკაში
ინტეგრაციის, splitting-ის, წაშლისა და ინდექსის ტესტები (მათ შორის
კონფლიქტური სცენარები), გეითვეიში ოთახის, გადაგზავნისა და permission
ტესტები, ფრონტენდზე სინქრონიზაციის ლოგიკის ტესტები. ბინარულ პროტოკოლს
ორივე მხარეს protocol drift ტესტები იცავს. CI (GitHub Actions) ყოველ
ცვლილებაზე უშვებს ფორმატირების, ლინტერების, ტესტებისა და უსაფრთხოების
შემოწმებებს (\texttt{cargo audit}, \texttt{npm audit},
\texttt{govulncheck}).

\paragraph{დარჩენილი შეზღუდვები.}
აღსანიშნავია ისიც, რაც სისტემას ჯერ არ შეუძლია: სესია ერთ დოკუმენტზეა
შემოსაზღვრული (მრავალფაილიანი პროექტების მხარდაჭერა არ არის), მონაწილეთა
კურსორები ერთმანეთისთვის ჯერ არ ჩანს, ხოლო ჰოსტის გათიშვა სესიას
ასრულებს. ეს პუნქტები სამომავლო განვითარების ბუნებრივი კანდიდატებია
(იხ. დასკვნა).
